\emph{Setting.} By Proposition \claimref{synthesis.operator-join-ordering-constants} the oscillators obey $[Z^{\alpha}_{i}, \bar{Z}_{l\beta}] = \hbar\,\delta_{il}\delta^{\alpha}_{\beta}$, $\hat{B}$ is the bracket $B_{\Lambda} = \br{\pi^{1}\pi^{2}} + \lambda[\omega^{1}\omega^{2}]$ as a polynomial in the $Z$ alone, $\hat{B}^{\dagger}$ its conjugate, a polynomial in the $\bar{Z}$ alone, $\hat{h}_{il} = \tfrac12(\bar{Z}_{i}\cdot Z^{l} + Z^{l}\cdot\bar{Z}_{i})$, $\hat{C} = \tfrac12\sum_{i,l}\hat{h}_{il}\hat{h}_{li} - \tfrac14(\tr\hat{h})^{2}$, and the palatial form of $\hat{m}^{2}$ is that of clause (3) there. By Fact \claimref{mass.imported-palatial-quadratics-homogeneity-charge}, in the realisation $\bar{Z} = -\hbar\partial$ the symmetrised identity quadratic of one constituent acts on functions of homogeneity $n$ as $-(n + 2)\hbar$; this is $\hat{h}_{ii}$ for constituent $i$.

\emph{(1) and (2).} Write $\mathrm{ad}(\hat{h}_{il})$ for the commutator with $\hat{h}_{il}$. From the canonical relations, $[\bar{Z}_{i}\cdot Z^{l}, Z^{\alpha}_{k}] = -\hbar\,\delta_{ik}Z^{\alpha}_{l}$ and $[\bar{Z}_{i}\cdot Z^{l}, \bar{Z}_{k\alpha}] = \hbar\,\delta_{lk}\bar{Z}_{i\alpha}$, and the symmetrisation changes $\hat{h}_{il}$ by a constant only for $i = l$, so $\mathrm{ad}(\hat{h}_{il})$ acts on polynomials in $Z$ and $\bar{Z}$ as a derivation which, on generators, moves the constituent label $i$ to $l$ (on $Z$) and $l$ to $i$ (on $\bar{Z}$). The bracket $\br{\pi^{1}\pi^{2}} = \epsilon^{A'B'}\pi^{1}_{A'}\pi^{2}_{B'}$ and $[\omega^{1}\omega^{2}] = \epsilon_{AB}\omega^{1A}\omega^{2B}$ are antisymmetric in the two constituents and bilinear, one factor from each: so $\mathrm{ad}(\hat{h}_{12})$ and $\mathrm{ad}(\hat{h}_{21})$ send them to the same form with two equal constituent labels, which vanishes by antisymmetry of $\epsilon$, and $\mathrm{ad}(\hat{h}_{ii})$ multiplies them by the same scalar, minus $\hbar$ in the normalisation of the Fact (each raises the homogeneity in $Z_{i}$ by one). Their conjugates, polynomials in $\bar{Z}$ of the same shape, get the opposite scalar and also commute with $\hat{h}_{12}$, $\hat{h}_{21}$. This is the infinitesimal form of the transformation law of Fact \claimref{families.imported-face-valued-bracket-quantum-numbers}: the bracket transforms by $\det(U)$ under the internal $U(2)$. Consequently $\hat{B}$ and $\hat{B}^{\dagger}$ have opposite charges and $\hat{B}^{\dagger}\hat{B}$, $\hat{B}\hat{B}^{\dagger}$ commute with every $\hat{h}_{il}$; $\hat{C}$ is the quadratic Casimir of the $\mathfrak{u}(2)$ spanned by the $\hat{h}_{il}$ (which close under commutators, by the derivation property) minus a multiple of the square of its central element, so it commutes with them; constants commute. This proves (1). The source term is the middle term of the expansion in Proposition \claimref{families.source-term-linear-tilt-mass-functional}(1), $\tfrac{2\Lambda}{3}\mathrm{Re}(\bar{b}\Omega)$, with $\Omega$ replaced by its vacuum value $A$, as that proposition's reading (2) prescribes ($b$ and the vacuum value of $\Omega$ independent), quantised as $\tfrac{\Lambda}{3}(\bar{A}\hat{b} + A\hat{b}^{\dagger})$; by (2) each term shifts both helicity eigenvalues by the same $\mp\hbar$.

The check realises all of this exactly on polynomials in eight variables with $\hbar = 1$ and $\Lambda = 3$: it verifies the stored commutator $[\hat{B}, \hat{B}^{\dagger}] = \hbar\lambda(\hat{h}_{11} + \hat{h}_{22})$ (which fixes the conjugation convention to be the store's), $\hat{h}_{ii} = -(n_{i} + 2)$ on each monomial, $[\hat{h}_{il}, \hat{m}^{2}] = 0$ for all four $(i, l)$, the charges of $\hat{b}$, $\hat{\Omega}$ and their conjugates, and that $\hat{S}$ sends every monomial of bidegree $(n_{1}, n_{2})$ into bidegrees $(n_{1} \pm 1, n_{2} \pm 1)$, on all monomials of total degree at most three; and that for a generic polynomial of bidegree $(2, 1)$ the image under $\hat{S}$ has no component of bidegree $(2, 1)$ while $\hat{m}^{2}$ keeps it there.

\emph{(3).} If $\hat{h}_{ii}$ is self-adjoint for a Hermitian form, its eigenspaces for distinct eigenvalues are orthogonal. For $\psi$ in the joint eigenspace $(t_{1}, t_{2})$, $\hat{S}\psi$, $\hat{b}\psi$ and $\hat{\Omega}\psi$ lie in the eigenspaces $(t_{1} \mp \hbar, t_{2} \mp \hbar)$, orthogonal to $\psi$, so the three expectations vanish. On a module on which both $\hat{h}_{ii}$ are scalars every vector is such a $\psi$. An irreducible positive-energy helicity module of one constituent is a module of the $\mathfrak{su}(2,2)$ of symmetrised traceless quadratics, which commute with that constituent's $\hat{h}_{ii}$; on the helicity modules used in the store $\hat{h}_{ii}$ is the scalar $2s\hbar$ (Proposition \claimref{synthesis.operator-join-casimir-shifts}(ii), via the Fact), so on a product of two such modules both are scalars. By (1), $\hat{m}^{2}$ preserves each joint eigenspace, so its eigenvectors there, including any with eigenvalue zero, are covered. A state $\psi = \sum_{t}\psi_{t}$ over joint eigenspaces has $\langle\psi, \hat{S}\psi\rangle = \sum_{t}\bigl(\langle\psi_{t - \hbar(1,1)}, \hat{S}\psi_{t}\rangle + \langle\psi_{t + \hbar(1,1)}, \hat{S}\psi_{t}\rangle\bigr)$, which needs two components whose helicity pairs differ by $\hbar(1,1)$, that is constituent helicities each differing by $\tfrac12$; such a state is not invariant under $\exp(I_{4}\phi\,\tr\hat{h}/\hbar)$, the determinant $U(1)$ (on the palatial Fock space numbers and the module's complex structure are read in $\mathrm{span}\{1, I_{4}\}$, the module read on the tower's floors, where $I_{4}$ commutes with the tower vectors), which commutes with $\hat{m}^{2}$ by (1). Because $\hat{m}^{2}$ preserves each joint eigenspace, its zero set is the direct sum of its zero sets in the separate eigenspaces, so a massless $\psi$ with $\langle\psi, \hat{S}\psi\rangle \ne 0$ exists only if $\hat{m}^{2}$ has zero modes in two eigenspaces whose eigenvalues differ by $\hbar(1,1)$ and $\hat{S}$ connects them. This is the ordinary mechanism of spontaneous breaking for a symmetric operator, and it is not examined here: the zero modes of $\hat{m}^{2}$ sector by sector are not computed, so the statement does not decide whether the massless locus of a direct sum of sectors supports the seed.

\emph{Scope.} The statement is about the source as the families' reading defines it, with the $\omega$-area a constant; the full cross term $\tfrac{\Lambda}{3}(\hat{b}^{\dagger}\hat{\Omega} + \hat{\Omega}^{\dagger}\hat{b})$, with both factors operators, has charge zero, but it is part of $\hat{m}^{2}$ itself and does not tilt anything. Classically, Theorem \claimref{synthesis.massless-composites-de-sitter-conformal-boundary} found the source's derivative at $b = 0$ vanishing on every massless composite with $b = 0$; (3) is the quantum counterpart, and stronger in one respect: on fixed-helicity modules the source vanishes on every state, massive ones included; weaker in another, since superpositions across sectors are not examined.

\emph{Reading (argued, at the observation level of the Fact).} By Fact \claimref{families.imported-face-valued-bracket-quantum-numbers} the bracket transforms under the global form as $(1,1)$ with $Y = +1$. The state the seed needs therefore breaks the $U(1)$ under which the bracket carries that hypercharge, which is the reservation the Fact records about what the bracket is an order parameter of; whether the join's massless locus supplies such a state is the open sub-question of (3); if it does not, the families' condensation programme must.
